% META: {"id": "1NSI-TC-EX-017-CDP", "chapitre": "1NSI-TYPES-CONSTRUITS", "type_objet": "coup_de_pouce", "exercice_ref": "1NSI-TC-EX-017", "status": "needs_review"}
\coupDePouce{1}{On te demande de prévoir les affichages de deux constructions par compréhension de liste, puis de réécrire l'une d'elles avec une boucle \lstinline{for} classique.}
\coupDePouce{2}{La syntaxe \lstinline{[n for n in notes if n >= 10]} signifie : « pour chaque \lstinline{n} dans \lstinline{notes}, garde \lstinline{n} seulement si \lstinline{n >= 10} ». Pour réécrire, utilise un tableau vide, une boucle et \lstinline{append}.}
